% ============================================================================
% Living spec of the MILP implemented in methods/manual/jobs/milp_jobs_v2_gurobipy.py
% (partition-based, job-level blocking semantics).  Kept in sync with the code;
% the polished journal version lives in papers/jobs_extension/milp.tex.
%
% Last synced: 2026-07-27, against the relaxed Mode-A model (commit 0fb58ea:
% closed before/after bounds, no eta clearance).
% ============================================================================

\section{MILP}

Every access of a blocked aircraft is explicitly \emph{classified}, by a
partition constraint, into one of the four mutually exclusive states the
problem allows: before the front aircraft's stay, after it, in an
inter-job gap (Mode~B), or strictly inside an interruptible job
(Mode~C).  An access that fits none of the four states --- an attempt to
enter or leave a position whose blocker is mid-way through a
non-interruptible job --- makes the partition constraint unsatisfiable
and is therefore correctly infeasible.

\subsection{Sets}

\begin{tabular}{l c l}
    $\mathcal{R}$ & : & Aircraft \\
    $\mathcal{P}$ & : & Positions \\
    $\mathcal{J}_r = (j^r_1, \ldots, j^r_{N_r})$ & : & Ordered job chain of aircraft $r$ (linearised from the precedence list) \\
    $\mathcal{J} = \bigcup_{r} \mathcal{J}_r$ & : & Global job set \\
    $\mathcal{A} \subseteq \mathcal{P} \times \mathcal{P}$ & : & Blocking arcs: if $(p,p') \in \mathcal{A}$, position $p$ (front) blocks $p'$ (rear)
\end{tabular}

\subsection{Parameters}

\begin{longtable}{>{\hspace{0pt}}l c p{11.2cm}}
    $D_j$ & : & Nominal duration of job $j$ \\
    $I_j \in \{0,1\}$ & : & Interruptibility flag of job $j$ ($I_j = 1$ iff $j$ can be paused) \\
    $E_r$ & : & Earliest start time of aircraft $r$ \\
    $L_r$ & : & Target finish time (due date) of aircraft $r$ \\
    $\varepsilon$ & : & Tow time; minimum separation between consecutive aircraft at the same position \\
    $\mu$ & : & Gap time consumed by one Mode~B manoeuvre (cumulative per gap) \\
    $\delta$ & : & Working-time extension added to an interruptible job per Mode~C pause \\
    $\eta$ & : & Time granularity enforcing the strict interior of Mode~C (one day on the benchmark) \\
    $K^{\max}$ & : & Bound on interruptions per job; the code uses $K^{\max} = \max(2,\, 2(|\mathcal{R}| - 1))$ \\
    $H^{UB}$ & : & Upper bound on every time variable (see below) \\
    $M$ & : & Big-$M$ constant, $M = H^{UB} + \max(\varepsilon, \mu, \delta, \eta) + 1$ \\
    $W^{M}, W^{D}, W^{S}$ & : & Weights for makespan, delays and movements \\
\end{longtable}

The horizon bound implemented in \texttt{prepare\_data} is the
worst-case sequential schedule plus the interruption-extension budget:
\[
  H^{UB} \;=\; \max_r E_r \;+\; \sum_{j \in \mathcal{J}} D_j
  \;+\; \varepsilon \,\max(0,\, |\mathcal{R}| - 1)
  \;+\; 2\delta\, |\mathcal{R}|
  \;+\; 4 \max(\mu, \eta) \max(1,\, |\mathcal{A}|).
\]
Defaults when the instance carries no paper-\#2 fields:
$\mu = 1$, $\delta = 2$, $\eta = 1$.

\subsection{Variables}

\paragraph{Assignment and timing}

\begin{longtable}{>{\hspace{0pt}}l c p{11.2cm}}
    $x_{rp} \in \{0,1\}$ & : & 1 if aircraft $r$ is assigned to position $p$ \\
    $q_{rr'p} \in \{0,1\}$ & : & 1 if aircraft $r$ precedes $r'$ at position $p$ (both assigned there) \\
    $s_j,\, f_j \in [0, H^{UB}]$ & : & Start and finish time of job $j$ \\
    $\kappa_j \in \{0, \ldots, K^{\max}\}$ & : & Interruption counter of job $j$ (integer; fixed to 0 if $I_j = 0$) \\
    $s_r,\, f_r$ & : & Aircraft shortcuts: $s_r \equiv s_{j^r_1}$, $f_r \equiv f_{j^r_{N_r}}$ (aliases of the chain endpoints, not separate variables) \\
\end{longtable}

\paragraph{Blocking}

For every ordered aircraft pair $r \ne r'$, every arc
$(p,p') \in \mathcal{A}$ and every access side
$a \in \{\mathrm{in}, \mathrm{out}\}$, with the convention
$\tau^a_{r'} = s_{r'}$ if $a = \mathrm{in}$ and $\tau^a_{r'} = f_{r'}$
if $a = \mathrm{out}$:

\begin{longtable}{>{\hspace{0pt}}l c p{10.4cm}}
    $y_{rr'pp'} \in \{0,1\}$ & : & 1 iff $r$ is the front aircraft at $p$ and $r'$ the rear at $p'$ \\
    $z^{-,a}_{rr'pp'} \in \{0,1\}$ & : & 1 iff access $\tau^a_{r'}$ falls at or before the front's start (Mode~A) \\
    $z^{+,a}_{rr'pp'} \in \{0,1\}$ & : & 1 iff access $\tau^a_{r'}$ falls at or after the front's finish (Mode~A) \\
    $b^{Ba}_{rr'pp'g} \in \{0,1\}$ & : & 1 iff the access falls in inter-job gap $g$ of $r$ (Mode~B), $g = 1, \ldots, N_r - 1$ \\
    $b^{Ca}_{rr'pp'j} \in \{0,1\}$ & : & 1 iff the access falls strictly inside job $j$ of $r$ (Mode~C); fixed to 0 if $I_j = 0$ \\
\end{longtable}

\paragraph{Metrics}

\begin{longtable}{>{\hspace{0pt}}l c p{11.2cm}}
    $v^{D}_r \ge 0$ & : & Positive delay of aircraft $r$ past $L_r$ \\
    $m \ge 0$ & : & Makespan \\
    $n \ge 0$ & : & Total movement count (two per manoeuvre) \\
\end{longtable}

\subsection{Constraints}

\paragraph{Assignment, duration, chain, earliest start}

\begin{align}
\sum_{p \in \mathcal{P}} x_{rp} &= 1
&& \forall r \in \mathcal{R} \\
%
f_j &= s_j + D_j + \delta\, \kappa_j
&& \forall j \in \mathcal{J} \\
%
\kappa_j &= 0
&& \forall j \in \mathcal{J} \;:\; I_j = 0 \\
%
s_{j^r_{g+1}} &\ge f_{j^r_g}
&& \forall r \in \mathcal{R},\; g = 1, \ldots, N_r - 1 \\
%
s_r &\ge E_r
&& \forall r \in \mathcal{R}
\end{align}

\paragraph{Same-position sequencing}

For every unordered pair $r < r'$ and every $p \in \mathcal{P}$ the code
creates both directed binaries and links them to the assignment:

\begin{align}
q_{rr'p} + q_{r'rp} &\ge x_{rp} + x_{r'p} - 1 \\
q_{rr'p} + q_{r'rp} &\le 1 \\
q_{rr'p} &\le x_{rp}, \quad q_{rr'p} \le x_{r'p}
\qquad\text{(and symmetrically for } q_{r'rp}\text{)} \\
%
s_{r'} &\ge f_r + \varepsilon - M (1 - q_{rr'p}) \\
s_{r} &\ge f_{r'} + \varepsilon - M (1 - q_{r'rp})
\end{align}

\paragraph{Placement compatibility}

For every $r \ne r'$ and $(p,p') \in \mathcal{A}$, $y_{rr'pp'}$ is the
conjunction of the two placements:

\begin{align}
y_{rr'pp'} \le x_{rp}, \qquad
y_{rr'pp'} \le x_{r'p'}, \qquad
y_{rr'pp'} \ge x_{rp} + x_{r'p'} - 1
\end{align}

\paragraph{Exhaustive partition of every access}

For every $r \ne r'$, $(p,p') \in \mathcal{A}$,
$a \in \{\mathrm{in}, \mathrm{out}\}$:

\begin{align}
z^{-,a}_{rr'pp'} + z^{+,a}_{rr'pp'}
 + \sum_{g=1}^{N_r-1} b^{Ba}_{rr'pp'g}
 + \sum_{j \in \mathcal{J}_r} b^{Ca}_{rr'pp'j}
 \;=\; y_{rr'pp'}
\label{eq:doc-partition}
\end{align}

When the blocking interaction is materialised ($y = 1$) the access must
take exactly one of the four states; when it is not ($y = 0$) no
indicator may fire.

\paragraph{Mode A --- before / after the front stay}

\begin{align}
\tau^a_{r'} &\le s_r + M\bigl(1 - z^{-,a}_{rr'pp'}\bigr)
\label{eq:doc-zminus} \\
\tau^a_{r'} &\ge f_r - M\bigl(1 - z^{+,a}_{rr'pp'}\bigr)
\label{eq:doc-zplus}
\end{align}

\paragraph{Mode B --- inter-job gap}

For every gap $g$ of the front aircraft $r$, with $j^r_g$ and
$j^r_{g+1}$ the jobs around it:

\begin{align}
\tau^a_{r'} &\ge f_{j^r_g} - M\bigl(1 - b^{Ba}_{rr'pp'g}\bigr) \\
\tau^a_{r'} &\le s_{j^r_{g+1}} + M\bigl(1 - b^{Ba}_{rr'pp'g}\bigr)
\end{align}

Each Mode~B event consumes $\mu$ units of its gap, cumulatively over
all rear neighbours and both access sides:

\begin{align}
s_{j^r_{g+1}} - f_{j^r_g} \;\ge\;
\mu \!\!\sum_{\substack{r' \ne r \\ (p,p') \in \mathcal{A} \\ a}}\!\! b^{Ba}_{rr'pp'g}
&& \forall r \in \mathcal{R},\; g = 1, \ldots, N_r - 1
\end{align}

\paragraph{Mode C --- strict interior of an interruptible job}

\begin{align}
\tau^a_{r'} &\ge s_j + \eta - M\bigl(1 - b^{Ca}_{rr'pp'j}\bigr) \\
\tau^a_{r'} &\le s_j + D_j + \delta\bigl(\kappa_j - b^{Ca}_{rr'pp'j}\bigr) - \eta
              + M\bigl(1 - b^{Ca}_{rr'pp'j}\bigr)
\label{eq:doc-bc-ub}
\end{align}

The upper bound~\eqref{eq:doc-bc-ub} excludes the access's own
contribution ($-\delta\, b^{Ca}$) from the extended finish, preventing
the self-justifying interruption in which an access becomes feasible
only because the extension it itself triggers pushes the job past it.

\paragraph{Interruption counter and movement count}

\begin{align}
\kappa_j &= \!\!\sum_{\substack{r' \ne r \\ (p,p') \in \mathcal{A} \\ a}}\!\! b^{Ca}_{rr'pp'j}
&& \forall r \in \mathcal{R},\; j \in \mathcal{J}_r \\
%
n &= 2 \sum_{(p,p') \in \mathcal{A}}
     \sum_{\substack{r,r' \in \mathcal{R} \\ r \neq r'}}
     \sum_{a} \Bigl(
       \sum_{g=1}^{N_r-1} b^{Ba}_{rr'pp'g}
       + \sum_{j \in \mathcal{J}_r} b^{Ca}_{rr'pp'j}
     \Bigr)
\end{align}

\paragraph{Delays and makespan}

\begin{align}
v^{D}_r &\ge f_r - L_r
&& \forall r \in \mathcal{R} \\
m &\ge f_r
&& \forall r \in \mathcal{R}
\end{align}

\subsection{Objective Function}

\begin{align}
\min \quad W^{M} m + W^{D} \sum_{r \in \mathcal{R}} v^{D}_{r} + W^{S} n
\end{align}

\begin{remark}[Relaxed Mode-A bounds]
The before/after constraints
\eqref{eq:doc-zminus}--\eqref{eq:doc-zplus} use \emph{closed} bounds
$\tau \le s_r$ / $\tau \ge f_r$: Mode~A only requires the front
position to be vacant at the access instant, so any access at or
outside the stay endpoints $[s_r, f_r]$ is admissible, matching the
problem statement and the compliance checker
(\texttt{problems/jobs/checker.py}).  An earlier version of the model
required a full $\eta$ clearance ($\tau \le s_r - \eta$ /
$\tau \ge f_r + \eta$), which was strictly more conservative than the
declared problem and cut off feasible schedules at $\pm\varepsilon$ of
a front stay; the forensic audit of 2026-07-23 (commit
\texttt{96d2595}, e.g.\ \texttt{two\_rows\_loose} seed~8 under the
delay-priority profile) exposed it and commit \texttt{0fb58ea} relaxed
the bounds.  The MILP cache was invalidated and the full 870-run
battery was re-executed on 2026-07-23/24 under the relaxed model.
\end{remark}

\begin{remark}[Strict interior via $\eta$]
$\eta$ is an instance datum (one day on the benchmark grid), not a
solver tolerance; it appears only in the Mode~C interior bounds, whose
margins therefore sit far above Gurobi's feasibility tolerances.  On
finer time grids $\eta$ should scale with the grid.
\end{remark}
